\section{The exact archimedean orbital scalar}
\label{sec:archimedean-orbital-scalar-v54}

This chapter computes an orbital measure for the actual setting.  It
does not choose another Haar measure at infinity.  The conclusion is
the original \texttt{PushInf}, with the literal constant
\[
 K_X(u')=\operatorname{gaugeU}(X.B.I)\,
       \operatorname{volU}(k)^d\,T_w(u')^{-kd},
 \qquad
 \operatorname{volU}(k)=\prod_{j=1}^{k}\frac{(2\pi)^j}{(j-1)!}.
 \tag{AOS.1}
\]
Here $k=D.R.k$, $d=[F:\mathbb Q]$ is the number of real places of
$F=X.B.I.F$, and $T_w(u')$ is the source's positive rational integer
viewed in $\mathbb R$.  The proof explains separately the determinant,
the compact-unitary volume, and the character normalization.  Their
combination leaves no unspecified proportionality constant.

The compact-stabilizer argument also gives the original
\texttt{ExInf}, \texttt{HtotInf}, \texttt{WeilInf},
\texttt{AinfValue}, and \texttt{AprimeInf}.  These are archimedean
statements on a supplied setting.  Arithmetic quotient-volume
comparisons, the finite-place factors, and the ambient automorphic
family remain separate inputs to their later consumers.

\subsection{The supplied objects and the normalization to be computed}

Fix any original datum $D$, a Hida setting $X$, its actual adelic bridge
$h_b$, two sets $F_{dN},F_{dN'}$ in $X.B.I.G_\infty$, and words
$u,u'$.  The original hypothesis \texttt{HypX} says that $X$ has its
measure ties through this bridge, its level subgroup is normalized,
and the two specified sets are the required unipotent fundamental
domains.  In particular, it includes
\[
 \mu_{G,\infty}
   =(\operatorname{toInf}_{h_b})_*
           \left(\bigotimes_{v\mid\infty}\mu_{G,v}^{\rm old}\right),
 \qquad
 (\iota_H)_*\mu_{H,\infty}
    =(\operatorname{toInf}_{h_b})_*\mu_H(u').
 \tag{AOS.2}
\]
The first measures on the right are the original
\texttt{Formalism.REAL.muG}; the second equation is the source's
ambient-measure tie for the actual stabilizer.  We retain the other
parts of \texttt{HypX}, even when an individual calculation needs only
these two equations.

Put $E=X.B.I.E$, $r=D.R.rS$, $a=2r+1$ and $n=X.B.I.n=k+a$.
The rational ordered tuple under consideration is
\[
 Y_0=(e_1,\ldots,e_r,z_{u'},f_1,\ldots,f_r),
 \qquad T_0=Y_0^*JY_0.
 \tag{AOS.3}
\]
The order in (AOS.3) is precisely \texttt{tupH(frameOf X u')};
we have merely written its vectors as columns.  At infinity let
$V_\infty=\operatorname{Mat}_{n\times a}(E\otimes_F F_\infty)$,
$\mathcal H_\infty=\operatorname{Herm}_a(E\otimes_F F_\infty)$,
$U_\infty$ be the independent-column locus, and
$\gamma(Y)=Y^*JY$.  The actual orbit map is
\[
 \pi_{Y_0}:G_\infty\longrightarrow V_\infty,
       \qquad g\longmapsto gY_0;
 \tag{AOS.4}
\]
it is \texttt{orbM X u'} in the source.

Let $\psi$ be any character satisfying the original
\texttt{MeasCompat\,$\psi$} and \texttt{IsStdChar}.  Its infinite
factor is
\[
       \psi_\infty(t)=\exp\left(2\pi i\sum_{v\mid\infty}t_v\right).
 \tag{AOS.5}
\]
For the setting obtained from GMR19 this is the one character already
chosen before all words, places and auxiliary tuples.  We do not
choose a character again in any calculation below.

The measures $\mu_\infty,\nu_\infty$ are \emph{arbitrary} witnesses of
the source's \texttt{SdW} and \texttt{SdHerm}.  Similarly
$\omega_T$ is any family satisfying \texttt{IsFibreFamily} for these
measures, the Gram map and $U_\infty$.  Thus it has the support,
disintegration and continuity properties stated in (GMR.3), on the
whole Hermitian parameter space.  GMR2--GMR5 prove existence and
uniqueness at this exact scope, including the finite product of real
places.  They will identify arbitrary witnesses with the explicit
densities calculated below.

Finally write $J_v=\operatorname{diag}(\epsilon_{v,1},\ldots,
\epsilon_{v,n})$, with $\epsilon_{v,i}\in\{1,-1\}$, for the original
real model.  Its Haar measure is $\mu_{G,v}^{\rm old}$, and its
source gauge is $g_v=\operatorname{Gauge.gaugeU}(\epsilon_v)$.
The original antecedent of \texttt{PushInf} is
$\operatorname{GaugeDens}(\epsilon_v)$ for every $v$.
AOS5 also proves this universal gauge-density assertion directly;
in particular $g_v>0$.  The product in (AOS.1) is
$\operatorname{gaugeU}(X.B.I)=\prod_v g_v$.

\subsection{The two self-dual measures and the real Jacobian}

% BEGIN AOS1
\paragraph{AOS1. The exact infinite-place pairings and self-dual densities.}
\label{claim:AOS1}
Use the supplied embedding $\varphi_v:E\hookrightarrow\mathbb C$
and supplied matrix $P_v$ with
$P_v^*\varphi_v(J)P_v=J_v$.  The tuple coordinate is
$W_v=P_v^{-1}\theta_v(Y)$, column by column, where $\theta_v$ is
the actual completion of $\varphi_v$ to $E\otimes_F F_\infty$.
The Hermitian coordinate is $S_v=\theta_v(S)$, with no conjugation by
$P_v$.  These coordinates are real linear homeomorphisms onto
$\prod_v\operatorname{Mat}_{n\times a}(\mathbb C)$ and
$\prod_v\operatorname{Herm}_a(\mathbb C)$, respectively.  They carry
the source pairings to
\[
 \begin{split}
 B(W,Z)&=2\sum_v\operatorname{Re}\operatorname{tr}(W_v^*J_vZ_v),\\
 B_H(S,T)&=\sum_v\operatorname{Re}\operatorname{tr}(S_vT_v).
 \end{split}
 \tag{AOS.6}
\]
Consequently \emph{every} allowed self-dual pair becomes
\[
             d\mu_\infty=2^{nad}\prod_v dW_v,
       \qquad d\nu_\infty=\prod_v dT_v,
 \tag{AOS.7}
\]
where $dW_v$ is coordinate real-imaginary Lebesgue measure and $dT_v$
is Euclidean volume for the Frobenius inner product on the real
Hermitian subspace.

\begin{proof}
The CM algebra at a real place is $\mathbb C$, with involution complex
conjugation.  Its trace to $\mathbb R$ is $z+\overline z$.
One can verify this without a choice of complex coordinates: the two
extensions of the real embedding to $E$ are $\varphi_v$ and
$\overline{\varphi_v}$, and their sum is the trace.  Base change
preserves the trace.  The completed embedding is bijective because a
quadratic CM basis gives two real linearly independent complex
numbers.  Applying it at every real place gives the asserted finite
product homeomorphism.  Multiplication by the supplied invertible
$P_v^{-1}$ then gives the tuple coordinates.

The pairing \texttt{BR} is the sum over columns of the quadratic
algebra trace of $x^*Jy$.  Substituting $x=P_vw$, $y=P_vz$ gives the
first equation of (AOS.6).  The definition of \texttt{pairHerm} has a
factor $1/2$ in front of the quadratic trace.  Since
$\operatorname{tr}(ST)$ is real for Hermitian $S,T$, this gives the
second equation of (AOS.6), with no factor two remaining.

In an orthonormal real basis of tuple space, the first form has
diagonal coefficients $2\epsilon_{v,i}$, each repeated for the real
and imaginary coordinate of each column.  Its absolute determinant
is $2^{2nad}$.  GMR2's Fourier-inversion calculation therefore gives
self-dual density $2^{nad}$.  On the Hermitian space the second form
is precisely the Frobenius inner product, so its determinant in an
orthonormal basis is one.  In particular its self-dual measure is
Frobenius volume.  If one instead uses unscaled off-diagonal real and
imaginary entries as coordinates, this latter volume has density
$2^{a(a-1)/2}$; that factor is already contained in our use of an
orthonormal Hermitian basis and must not be inserted a second time.

Existence follows from Fourier inversion and finite products.  GMR2
gives uniqueness even for the source's conditional Fourier-inversion
predicate, so (AOS.7) holds for every original witness.  This argument
also shows why the determinant of $P_v$ does not create another
factor: the measures are transported together with their pairings.
\end{proof}
% END AOS1

% BEGIN AOS2
\paragraph{AOS2. Two determinant identities for complex matrices.}
\label{claim:AOS2}
For $B\in\operatorname{GL}_a(\mathbb C)$, right multiplication by $B$
on $\operatorname{Mat}_{b\times a}(\mathbb C)$ has absolute real
determinant $|\det B|^{2b}$.  The real linear map
\[
       H\longmapsto B^{-*}HB^{-1}
       \quad\hbox{on }\operatorname{Herm}_a(\mathbb C)
\]
has absolute determinant $|\det B|^{-2a}$, with Frobenius volume on
both sides.  The assertions include zero dimensions.

\begin{proof}
The real determinant of a complex linear map is its complex
determinant times its conjugate.  Apply this to each of the $b$ rows
to obtain the first formula.  For the second, use a singular-value
decomposition $B=U\operatorname{diag}(s_1,\ldots,s_a)V$ with $s_i>0$.
Unitary conjugations are Frobenius isometries, so they have absolute
real determinant one.  On the diagonal Hermitian coordinate indexed
by $i$, conjugation by $\operatorname{diag}(s_i^{-1})$ multiplies by
$s_i^{-2}$.  On each two-dimensional real off-diagonal coordinate
indexed by $i<j$, it multiplies both coordinates by $(s_is_j)^{-1}$.
Their product is
\[
       \prod_i s_i^{-2}\prod_{i<j}(s_is_j)^{-2}
          =\prod_i s_i^{-2a}=|\det B|^{-2a}.
 \tag{AOS.8}
\]
When $a=0$ or $b=0$, the relevant maps are maps of the zero vector
space and the empty determinant and product are one, as required.
\end{proof}
% END AOS2

% BEGIN AOS3
\paragraph{AOS3. The full original \texttt{PlaceJac} determinant.}
\label{claim:AOS3}
Let $J=\operatorname{diag}(\epsilon_1,\ldots,\epsilon_n)$ with signs
$\epsilon_i\in\{1,-1\}$.  Let $S$ be any set of $k$ coordinate
indices, put $n=a+k$, and let $w\in\operatorname{Mat}_{n\times a}
(\mathbb C)$ vanish in the rows indexed by $S$.  Suppose
$T=w^*Jw$ is invertible.  Let $\mathfrak m$ consist of the matrices
$Z$ satisfying $Z^*J+JZ=0$ whose $S\times S$ block is zero, with its
Frobenius metric.  Then
\[
 \mathcal D_w:\mathfrak m\oplus\operatorname{Herm}_a(\mathbb C)
       \longrightarrow\operatorname{Mat}_{n\times a}(\mathbb C),
 \qquad
 \mathcal D_w(Z,H)=Zw+\tfrac12wT^{-1}H
 \tag{AOS.9}
\]
is a real linear isomorphism, and
\[
        |\det_{\mathbb R}\mathcal D_w|
                =2^{-na}|\det T|^k.
 \tag{AOS.10}
\]
Equivalently, its pushforward of product Euclidean volume is
$2^{na}|\det T|^{-k}$ times tuple volume.  This is the full original
\texttt{PlaceJac}, for arbitrary finite orthonormal coordinate
spaces on its two factors.  No positivity assumption on the signs
indexed by $S$ is needed for this determinant statement.

\begin{proof}
Permute coordinates so that $S$ comes first.  These permutations are
Euclidean isometries.  Write
\[
       J=\begin{pmatrix}D_0&0\\0&D_1\end{pmatrix},
       \qquad w=\begin{pmatrix}0\\B\end{pmatrix}.
\]
Here $D_0,D_1$ are diagonal sign matrices and $B$ is invertible,
because $T=B^*D_1B$ is invertible.  The defining Lie algebra equations
give every $Z\in\mathfrak m$ uniquely as
\[
 Z=\begin{pmatrix}
       0&A\\-D_1A^*D_0&iD_1Q
    \end{pmatrix},\qquad Q^*=Q.
 \tag{AOS.11}
\]
Since the sign matrices are Euclidean unitary,
$\|Z\|_F^2=2\|A\|_F^2+\|Q\|_F^2$.
Thus $(R,Q)\mapsto Z$ with $A=R/\sqrt2$ is a real isometry
from coordinate $k\times a$ complex matrices and Frobenius Hermitian
matrices onto $\mathfrak m$.

In these coordinates (AOS.9) has upper block
$RB/\sqrt2$ and lower block
\[
       D_1\left(iQB+\tfrac12B^{-*}H\right)
       =D_1\left(iQ+\tfrac12B^{-*}HB^{-1}\right)B.
 \tag{AOS.12}
\]
The upper-block determinant is $2^{-ka}|\det B|^{2k}$ by AOS2.
For the lower block, the Hermitian congruence in $H$ contributes
$|\det B|^{-2a}$; multiplication by $1/2$ on its real $a^2$
coordinates contributes $2^{-a^2}$.  The map
$(Q,H')\mapsto iQ+H'$ is an isometry from two orthogonal Hermitian
spaces onto all complex matrices.  Finally right multiplication by
$B$ contributes $|\det B|^{2a}$, and left multiplication by $D_1$
contributes one.  The lower-block determinant is therefore
$2^{-a^2}$.  Multiplying the blocks gives
$2^{-a(k+a)}|\det B|^{2k}$, which is (AOS.10) because
$|\det T|=|\det B|^2$.

Each map used is invertible, which also proves that $\mathcal D_w$
is an isomorphism.  Any coordinates $L,H$ allowed in the original
\texttt{PlaceJac} are onto Frobenius isometries by its exact norm and
surjectivity hypotheses.  Composing with them changes no absolute
determinant.  The usual linear change-of-variables formula takes the
reciprocal of (AOS.10), exactly as in the original measure identity.
The calculation uses empty products when $a=0$ or $k=0$, so both
boundary cases are included.
\end{proof}
% END AOS3

% BEGIN AOS4
\paragraph{AOS4. A Gram chart and the exact density of its fiber.}
\label{claim:AOS4}
Keep $J,w,T$ as in AOS3.  Near $T$ there is a real analytic matrix
function $b(Q)$ such that
\[
 b(T)=1,\qquad b(Q)^*Tb(Q)=Q,
       \qquad Db_T(H)=\tfrac12T^{-1}H.
 \tag{AOS.13}
\]
For sufficiently small neighborhoods of $0$ and $T$, the map
\[
       \Psi(Z,Q)=\exp(Z)wb(Q)
 \tag{AOS.14}
\]
is a diffeomorphism onto an open subset of the independent tuple
locus, and $\gamma(\Psi(Z,Q))=Q$.  If tuple and Hermitian measures
have the self-dual normalizations of AOS1, the fiber measure in the
chart $Z\mapsto\exp(Z)w$ has a continuous positive density $j_T(Z)$
and
\[
                     j_T(0)=|\det T|^k.
 \tag{AOS.15}
\]
This holds for every source fiber-family witness.

\begin{proof}
For $Q=T+H$ sufficiently close to $T$, the convergent matrix power
series for the square root defines
$b(Q)=(1+T^{-1}H)^{1/2}$, using the branch equal to one at $H=0$.
The matrix $A=T^{-1}H$ is $T$-self-adjoint: $A^*T=TA$.
Its real-coefficient power series therefore satisfies
$b(Q)^*T=Tb(Q)$.  Since $b(Q)^2=1+A$, (AOS.13) follows;
its derivative is the first-order term of the series.  This argument
does not assume that $T$ is positive definite.

The matrices $Z\in\mathfrak m$ are in the Lie algebra of $U(J)$,
so $\exp(Z)^*J\exp(Z)=J$.  Consequently (AOS.14) has Gram matrix
$Q$.  Its derivative at $(0,T)$ is exactly (AOS.9).  AOS3 and the
real inverse function theorem give the asserted local diffeomorphism.
Shrink to a product neighborhood on which its Jacobian does not
vanish and its tuples are independent.

Write $dZ$ for Euclidean volume on $\mathfrak m$ and $dQ$ for
Frobenius Hermitian volume.  The change-of-variables formula gives
\[
       d\mu(\Psi(Z,Q))
         =2^{na}|\det_{\mathbb R}D\Psi_{(Z,Q)}|\,dZ\,dQ.
 \tag{AOS.16}
\]
In these coordinates the Gram map is projection to $Q$.  GMR4's
fiber construction therefore has density
$2^{na}|\det D\Psi_{(Z,Q)}|\,dZ$ on each fiber in this chart.
GMR5 identifies any permitted fiber family with that construction
at \emph{every} $Q$, using continuity of compactly supported fiber
integrals; an almost-everywhere disintegration is not substituted
for that conclusion.  Evaluating (AOS.16) at $(0,T)$ and applying
(AOS.10) proves (AOS.15).  Positivity and continuity follow from
the nonzero continuous Jacobian.
\end{proof}
% END AOS4

\subsection{The prescribed Haar density and the compact volume}

% BEGIN AOS5
\paragraph{AOS5. The exact Haar density and the full original \texttt{GaugeDens}.}
\label{claim:AOS5}
For any sign matrix $J$, let $d_G$ be the left invariant Riemannian
volume on $U(J)$ whose inner product at the identity is Frobenius
inner product on
$\mathfrak g_J=\{Z:Z^*J+JZ=0\}$.  If $\mu$ is any Haar measure
on this group, its density relative to $d_G$ is the positive scalar
\[
       c_\mu=\lim_{\rho\downarrow0}
          \frac{\mu(\exp B_{\mathfrak g_J}(0,\rho))}
               {\operatorname{vol}_{\mathfrak g_J}
                    (B_{\mathfrak g_J}(0,\rho))}.
 \tag{AOS.17}
\]
For the source measure $\mu_G^{\rm old}$ this is precisely its
$\operatorname{gaugeU}$, so $\mu_G^{\rm old}=g\,d_G$.
In particular the full original $\operatorname{GaugeDens}(\epsilon)$
holds for every dimension and every sign function $\epsilon$.
The comparison measure $d_G$ is used
to calculate the density; it does not replace $\mu_G^{\rm old}$.

\begin{proof}
The Lie algebra appearing here is the source's \texttt{gLie}.
Indeed differentiating $\exp(tZ)^*J\exp(tZ)=J$ at zero gives
$Z^*J+JZ=0$.  Conversely this equation implies that the derivative
of $\exp(tZ)^*J\exp(tZ)$ is zero, so its value is $J$ for every
$t$.  Taking the real span in the source definition changes nothing.
The source realification map is a Frobenius isometry.

For completeness, smooth local group coordinates need no additional
structure theorem here.  The Cayley map
\[
       Z\longmapsto (1+Z/2)(1-Z/2)^{-1}
\]
is a chart from a neighborhood of zero in this Lie algebra to a
neighborhood of the identity in $U(J)$.  The equation
$Z^*J+JZ=0$ verifies unitarity by multiplication; its inverse is
$g\mapsto2(g-1)(g+1)^{-1}$ wherever the latter is defined.
Both maps are smooth, and the first derivative at zero is the
identity.  Translate this chart to obtain the group manifold.

The left-translated Frobenius inner product defines a smooth positive
volume density.  Its associated measure is locally finite, positive
on nonempty open sets, and left invariant.  It is therefore Haar.
Uniqueness of Haar measure gives $\mu=c_\mu d_G$ for a finite
$c_\mu>0$.  On a sufficiently small Lie algebra ball, exponential
is a diffeomorphism, and the pullback of $d_G$ is
$J_{\exp}(Z)\,dZ$, where $J_{\exp}$ is continuous and
$J_{\exp}(0)=1$ because $D\exp_0$ is the identity.  The ratio in
(AOS.17) is the average of $c_\mu J_{\exp}$ on this ball.  Uniform
continuity near zero makes that average tend to $c_\mu$.

The original \texttt{expSet} is this exponential image, and its
\texttt{ratioAt} uses exactly the Lie algebra volume just described.
Thus the limit exists and is $c_\mu>0$.  The source defines
$\operatorname{gaugeU}$ by \texttt{limUnder}, which equals any
existing limit.  Consequently $g=c_\mu$, and both the convergence
and strict positivity clauses of \texttt{GaugeDens} follow.  This
proves the normalization from the literal gauge definition.  In
dimension zero the exponential ball and the group are singletons;
their Lie volume is one and the same assertion holds directly.
\end{proof}
% END AOS5

% BEGIN AOS6
\paragraph{AOS6. The compact unitary group has the required exact volume.}
\label{claim:AOS6}
For $U(j)$ with the Frobenius normalization of AOS5,
\[
                  d_{U(j)}(U(j))
                      =\prod_{i=1}^{j}\frac{(2\pi)^i}{(i-1)!}.
 \tag{AOS.18}
\]
In particular this proves the original \texttt{GaugeVolU}: under its
stated \texttt{GaugeDens} antecedent, the source measure
$\operatorname{muGauge}$ has total mass $\operatorname{volU}(j)$.

\begin{proof}
The group $U(j)$ is closed and bounded in its matrix space, hence
compact.  The map $q_j:U(j)\to S^{2j-1}$ taking the last column is
onto, and the stabilizer of the last standard basis vector is the
actual block $U(j-1)$.  Local orthonormal completion gives smooth
sections.  At the identity its vertical Lie algebra is the upper
$\mathfrak u(j-1)$ block.  Its Frobenius orthogonal complement has
matrices
\[
       \begin{pmatrix}0&z\\-z^*&it\end{pmatrix},
       \qquad \|Z\|_F^2=2\|z\|^2+t^2.
 \tag{AOS.19}
\]
The last-column derivative is $(z,it)$, whose round-sphere tangent
norm is $\|z\|^2+t^2$.  Its absolute determinant on orthonormal
horizontal coordinates is $2^{-(j-1)}$, since $z$ has real
dimension $2j-2$ and $t$ has dimension one.

Disintegrate the invariant group volume over the last-column map.
The fiber volume is the unchanged Frobenius volume on $U(j-1)$:
left multiplication is an isometry and the block inclusion preserves
Frobenius inner products.  The quotient density is invariant under
the transitive $U(j)$ action on the sphere.  Formula (AOS.19), or
ordinary change of variables in a local section times a fiber chart,
shows that this quotient volume is $2^{j-1}$ times round-sphere
area.  The equality of these densities at one point gives it
everywhere by invariance.  Thus
\[
 d_{U(j)}(U(j))
      =d_{U(j-1)}(U(j-1))\,2^{j-1}
                        \operatorname{area}(S^{2j-1}).
 \tag{AOS.20}
\]
The area is $2\pi^j/(j-1)!$: integrate $e^{-\|x\|^2}$ on
$\mathbb R^{2j}$, first as a product of $2j$ one-dimensional
Gaussian integrals, giving $\pi^j$, and then in polar coordinates.
The radial integral is
$\tfrac12\int_0^\infty e^{-s}s^{j-1}ds=(j-1)!/2$, by repeated
integration by parts.  Equation (AOS.20) therefore has successive
ratio $(2\pi)^j/(j-1)!$.

For $j=0$ the group and its Lie algebra have one point and the
normalized volume is one.  Induction proves (AOS.18), including
that boundary case.  Finally AOS5 gives
$\operatorname{muGauge}=g^{-1}\mu_G^{\rm old}=d_{U(j)}$,
so the source formula follows with its exact product convention.
\end{proof}
% END AOS6

% BEGIN AOS7
\paragraph{AOS7. The local density of the actual orbit pushforward.}
\label{claim:AOS7}
Suppose now that $J=\operatorname{diag}(I_k,D_1)$, and
$w=\binom{0}{B}$ with $B$ invertible.  Its pointwise stabilizer is
the actual block subgroup $H=U(k)$.  Let $\mu_G=g\,d_G$, with
$g>0$.  In the local orbit chart $Z\mapsto\exp(Z)w$ from AOS4,
$\pi_{w*}\mu_G$ has continuous positive density $q(Z)$, with
\[
                         q(0)=g\operatorname{volU}(k).
 \tag{AOS.21}
\]
The orbit map is proper, its pushforward is Radon, and every nonzero
invariant Radon measure on the same orbit is a positive multiple of
this pushforward.

\begin{proof}
An element fixing $w$ fixes the whole last $a$-dimensional coordinate
space, since $B$ is invertible.  Unitarity then makes it preserve
the orthogonal first $k$-space.  Hence it is exactly
$\operatorname{diag}(h,1)$, $h\in U(k)$; the converse is immediate.
The Frobenius splitting
$\mathfrak g_J=\mathfrak m\oplus\mathfrak u(k)$ is orthogonal.
Exponential and the inverse function theorem give a local section of
$G/H$ of the form $Z\mapsto\exp(Z)$.  AOS4 identifies this quotient
chart with the indicated orbit chart.

Pull back $\mu_G$ by $(Z,h)\mapsto\exp(Z)h$, using $dZ$ and
Frobenius Haar $d_H$ on the compact group $H$.  The resulting density
is a smooth positive function $D(Z,h)$.  At $(0,h)$, left translate
the derivative back by $h^{-1}$.  Its horizontal part is
$\operatorname{Ad}(h^{-1})Z$ and its vertical part is the usual
left-translated tangent vector in $\mathfrak h$.  Since $h$ is
Euclidean unitary and block diagonal, its adjoint action preserves
the Frobenius metric and the orthogonal splitting.  The derivative
therefore has absolute determinant one.  AOS5 implies
$D(0,h)=g$ for every $h\in H$.

This product map is a diffeomorphism onto the inverse image of a
sufficiently small quotient neighborhood: uniqueness of the local
section makes the $Z$ coordinate unique and then makes $h$ unique.
It follows that $q(Z)=\int_H D(Z,h)\,d_H(h)$.  Compactness gives
continuity and positivity, and AOS6 gives (AOS.21).

The same local-section description proves properness.  Cover a
compact subset of the orbit by finitely many smaller section
neighborhoods with compact closures inside section domains.  The
preimage of each corresponding compact piece lies in the product
of the compact section image with the compact $H$.  Their union is
compact.  Thus the orbit pushforward of the Radon Haar measure is
Radon.

For uniqueness, let $\rho$ be an invariant Radon measure on the
orbit and let $dh$ be probability Haar on $H$.  Integrating
$\int_H f(s(y)h)\,dh$ against $\rho(y)$ defines a left invariant
Radon measure on $G$, by the local-section argument of GMR11.
If $\rho\ne0$, invariance and transitivity make every nonempty open
set positive: otherwise its countably many translates would cover
the orbit with zero measure.  The lifted measure is therefore a
nonzero Haar measure, hence a finite positive multiple of $\mu_G$.
Pushing it back to the orbit recovers $\rho$, because $dh(H)=1$.
This proves the asserted uniqueness.  It also includes the zero
measure with scalar zero when that case is needed.
\end{proof}
% END AOS7

% BEGIN AOS8
\paragraph{AOS8. Exact one-place orbital measure, with no unknown scalar.}
\label{claim:AOS8}
Let $w\in\operatorname{Mat}_{n\times a}(\mathbb C)$ have invertible
Gram matrix $T=w^*Jw$, and suppose its orthogonal complement is
positive definite of dimension $k=n-a$.  Use the source character
$e^{2\pi it}$, its tuple and Hermitian self-dual measures, and any
permitted continuous fiber family $\omega$.  For the prescribed
Haar measure $\mu_G=g\,d_G$,
\[
            \pi_{w*}\mu_G
                 =g\operatorname{volU}(k)|\det T|^{-k}\,\omega_T.
 \tag{AOS.22}
\]

\begin{proof}
First suppose the orthogonal complement is a positive coordinate
$k$-space, so $w=\binom0B$.  Witt extension identifies its orbit
with the entire Gram fiber.  Since $T$ is invertible, that entire
fiber is inside the independent locus.  GMR4--GMR5 give an invariant
Radon fiber measure, positive on every nonempty relatively open set.
AOS7 therefore makes the orbit pushforward $c\omega_T$ for a finite
$c>0$.

Use the same local orbit coordinate $Z$ for both measures.  Its
fiber density is continuous with value $|\det T|^k$ by AOS4, and
its orbit-pushforward density is continuous with value
$g\operatorname{volU}(k)$ by AOS7.  Their proportionality forces
\[
       c=\frac{g\operatorname{volU}(k)}{|\det T|^k}.
\]
For example, integrate both densities over balls of radius $\rho$
and divide by their Euclidean volumes; as $\rho\downarrow0$ the
two averages tend to the stated positive values.  In a zero
dimensional chart the same calculation uses the singleton of
volume one.  This proves (AOS.22) without discarding a null set at
the base point.

For a general positive orthogonal complement, choose an orthonormal
basis of it.  Hermitian Gram--Schmidt on its nondegenerate
orthogonal complement extends this to an isometry from a coordinate
sign decomposition.  Equivalently, the elementary Witt extension
proved in GMR10 gives $c_0\in U(J)$ carrying a positive coordinate
$k$-space to this complement.  Thus $c_0^{-1}w$ has the preceding
block form, with the same Gram matrix $T$.

The unitary group is unimodular.  Here is the real version of
GMR14's trace-pairing argument.  The real Lie algebra
$\{Z:Z^*J+JZ=0\}$ is a real form of
$\operatorname{Mat}_n(\mathbb C)$: for the conjugate-linear
involution $\tau(Z)=-JZ^*J$, every complex matrix is the sum of
a fixed vector and $i$ times a fixed vector.  The pairing
$\operatorname{Re}\operatorname{tr}(AB)$ on this real form
complexifies to the nondegenerate matrix trace pairing.  It is
therefore nondegenerate and invariant under conjugation.  Thus
$(\det\operatorname{Ad}(g))^2=1$.  Conjugation multiplies a
left invariant volume density by the absolute determinant of its
derivative at the identity, which is consequently one.  Composing
conjugation with a left translation proves that every right
translation preserves every left Haar normalization.  Therefore
$\pi_w(g)=\pi_{c_0^{-1}w}(gc_0)$ gives identical pushforward
measures.  Applying the block calculation proves the general case.
The coordinate isometry is used only in this computation; the
measure on the original group and the actual orbit are unchanged.
\end{proof}
% END AOS8

\subsection{The source frame and all infinite-place witnesses}

% BEGIN AOS9
\paragraph{AOS9. The source frame has the prescribed determinant and compact stabilizer.}
\label{claim:AOS9}
For every original $D,X,u'$, the tuple (AOS.3) is independent,
has nondegenerate Hermitian Gram matrix, and at every real place
\[
          |\det\varphi_v(T_0)|=T_w(u').
 \tag{AOS.23}
\]
Its orthogonal complement is positive definite of dimension $k$.
The actual infinite pointwise stabilizer $H_\infty$ is a compact
group isomorphic to $\prod_{v\mid\infty}U(k)$.  In particular this
proves the original \texttt{HinfCompact} at its full scope; neither
measure compatibility nor a test function is required for compactness.

\begin{proof}
The setting's frame identities give, in the order $(e,z,f)$,
\[
       T_0=\begin{pmatrix}
           0&0&I_r\\0&-T_w(u')&0\\I_r&0&0
          \end{pmatrix}.
 \tag{AOS.24}
\]
These identities are transported by the supplied coordinate map and
field identification, rather than by a newly chosen frame.
Reordering as $(e,f,z)$ shows
$\det T_0=(-1)^{r+1}T_w(u')$.  The integer $T_w(u')$ is positive,
so (AOS.23) follows at every real embedding.  Invertibility of the
Gram matrix implies independence: a relation between the columns
would be in the kernel of $T_0$.

Each hyperbolic pair has one positive and one negative direction,
as is seen from $(e_i+f_i)/\sqrt2$ and $(e_i-f_i)/\sqrt2$.
The vector $z$ is negative, of squared length $-T_w(u')$.
Thus the tuple span has signature $(r,r+1)$.  The original setting
has exactly $r+1$ negative directions.  Its rank is
$n=k+m-1=k+2r+1$, because its $m$ is even and at least two,
and $r=m/2-1$.  Sylvester's law of inertia, applied to the orthogonal
direct sum with this nondegenerate tuple span, makes its complement
positive definite and of dimension $k$.

An isometry fixing each tuple vector fixes the tuple span pointwise
and acts arbitrarily unitarily on its orthogonal complement.
Conversely every such complementary unitary extends by the identity
on the span.  These two matrix constructions are inverse continuous
homomorphisms.  At each real place the stabilizer is therefore
isomorphic to $U(k)$; the completed CM coordinates give the finite
product assertion for the actual $H_\infty$.  Compactness follows
from compactness of $U(k)$.  The argument includes $r=0$ (one
negative vector) and the zero-dimensional complementary group.
\end{proof}
% END AOS9

% BEGIN AOS10
\paragraph{AOS10. Products preserve the exact fiber measure at every parameter.}
\label{claim:AOS10}
Under the completed coordinates of AOS1, let $\omega_{v,T_v}$ be
the local fiber measures.  For every $T=(T_v)$,
\[
                         \omega_T=\bigotimes_v\omega_{v,T_v}.
 \tag{AOS.25}
\]
This equality holds for every permitted global self-dual pair and
every permitted global fiber-family witness.  At nondegenerate $T$,
the product of the orbit maps pushes the prescribed product Haar
measure to the product of the local orbit measures.

\begin{proof}
The independent locus over the product algebra is the product of the
local independent loci: multiplying a linear relation by each
coordinate idempotent proves this equivalence in both directions.
The Gram map is componentwise.  On products
of the local submersion charts of GMR3--GMR4, the ambient Jacobian is
the product of their positive Jacobians.  The fiber density at a
fixed parameter is therefore their product, by ordinary Fubini.
These product charts cover the independent locus and their measures
agree on overlaps, so they construct the product fiber family.
The support and disintegration statements follow by Tonelli and a
countable chart cover.  A compactly supported continuous test is
covered by finitely many product charts; a partition of unity and
a common compact bound in each chart give continuity of its fiber
integral by dominated convergence.  This verifies the full original
fiber-family condition, including singular parameters.

AOS1 identifies every permitted global self-dual measure with the
corresponding product.  GMR5 now identifies any permitted global
fiber family with the product construction pointwise, not merely
almost everywhere.  This proves (AOS.25).  For orbit pushforwards,
equality on measurable rectangles is the definition of product
measure and product pushforward; a monotone-class argument gives
equality of measures.  Compact stabilizers make these orbit
pushforwards Radon by AOS7, so the needed product measures are
sigma finite.  No infinite restricted product occurs in this
archimedean assertion.
\end{proof}
% END AOS10

% BEGIN AOS11
\paragraph{AOS11. The full original \texttt{PushInf}.}
\label{claim:AOS11}
For every $D,X,h_b,F_{dN},F_{dN'},u,u',\psi$ satisfying the original
\texttt{HypX}, \texttt{MeasCompat\,$\psi$}, \texttt{IsStdChar}, and
all original \texttt{GaugeDens} antecedents, for every witness
$h_T$ of Hermitian symmetry of $T_0$, and for every original
self-dual pair and fiber family,
\[
             (\pi_{Y_0})_*\mu_{G,\infty}
                  =\operatorname{ofReal}(K_X(u'))\,\omega_{T_0}.
 \tag{AOS.26}
\]
Here $T_0$ on the right is the literal \texttt{hermOf} determined by
the rational Gram matrix and $h_T$.  Thus (AOS.26) is the original
\texttt{PushInf} equality, with its complete witness quantifiers.

\begin{proof}
The completed bridge intertwines the actual action with the product
of the real-model actions.  Explicitly,
$\theta_v(g)=P_vg_vP_v^{-1}$ and
$\theta_v(Y)=P_vW_v$, so
$P_v^{-1}\theta_v(gY)=g_vW_v$.  This uses the supplied bridge;
its equality on matrices determines the action without a further
isomorphism choice.  The first measure tie in (AOS.2) identifies
the prescribed $\mu_{G,\infty}$ with the pushforward of precisely
the product of the original real Haar measures.

AOS9 verifies the hypotheses of AOS8 at every real place.  AOS5
identifies the scalar of its original Haar measure with $g_v$.
Multiplying (AOS.22), using AOS10 and (AOS.23), gives the scalar
\[
       \prod_v g_v\operatorname{volU}(k)
                            |\det\varphi_v(T_0)|^{-k}
          =\left(\prod_vg_v\right)
               \operatorname{volU}(k)^dT_w(u')^{-kd}.
 \tag{AOS.27}
\]
It is exactly (AOS.1).  All its factors are strictly positive and
finite.  Transport back through the same coordinate homeomorphism
proves equality of the measures on the original tuple space.
The proof $h_T$ specifies an element of a fixed Hermitian subspace;
its underlying matrix is unchanged by a different symmetry proof.
Uniqueness in AOS1 and AOS10 proves the assertion for all the
self-dual and fiber witnesses quantified in the original statement.

No other part of the setting is altered.  In particular the actual
bridge, $\mu_{G,\infty}$, $\mu_{N,\infty}$ and
$\mu_{H,\infty}$ are the ones supplied before this argument.
The unused level and fundamental-domain conditions are retained
because they are antecedents of the literal original theorem.
\end{proof}
% END AOS11

% BEGIN AOS12
\paragraph{AOS12. A positive finite test for the full original \texttt{ConstTest}.}
\label{claim:AOS12}
For every datum satisfying the original \texttt{PushHyp}, there is a
measurable $\phi:V_\infty\to[0,\infty]$ such that
\[
 \begin{gathered}
             0<\int\phi\,d\omega_{T_0}<\infty,\\
 \int\phi\,d(\pi_{Y_0*}\mu_{G,\infty})
                  =K_X(u')\int\phi\,d\omega_{T_0}.
 \end{gathered}
 \tag{AOS.28}
\]
In fact $\phi$ may be taken continuous, nonnegative and compactly
supported, with support contained in the independent locus.

\begin{proof}
The independent locus is open and contains $Y_0$ by AOS9.  In any
fixed Euclidean coordinates, take $\delta>0$ such that the closed
ball of radius $2\delta$ about $Y_0$ is inside that locus.  Define
\[
             \phi(Y)=\max(1-\|Y-Y_0\|/\delta,0).
 \tag{AOS.29}
\]
Its support is compact and it is bounded by one.  The fiber measure
is Radon on the independent locus, so the integral is finite.
It is positive because $\phi\geq1/2$ on the relative fiber ball of
radius $\delta/2$, whose measure is positive by GMR4, or by the
positive local density in AOS4.  Regard (AOS.29) as an extended
nonnegative-valued function.  It is measurable and still has those
integral bounds.

The existential witnesses in \texttt{PushHyp} supply the bridge,
sets and first word needed for AOS11; its other clauses are exactly
the remaining antecedents and witness conditions of that theorem.
Apply (AOS.26) to the just constructed test.  This proves (AOS.28)
with the prescribed constant and no cancellation of a zero or
infinite integral.  It is the full original \texttt{ConstTest}.
\end{proof}
% END AOS12

\subsection{Existence, the actual compact Haar mass, and all sections}

% BEGIN AOS13
\paragraph{AOS13. The full original \texttt{ExInf}, with all witnesses.}
\label{claim:AOS13}
For every supplied setting, bridge, words, sets and character under
the antecedents of \texttt{ExInf}, the nested \texttt{AllSome} assertion
in that definition holds: tuple self-dual measures exist; for every
such measure Hermitian self-dual measures exist; and for every such
pair a full continuous fiber family exists.

\begin{proof}
Apply GMR2 at the finite product of real places with the source
pairings and character.  AOS1 gives their exact coordinate forms
and hence nondegeneracy, so it also supplies a direct realization
of the two self-dual measures.  GMR4 constructs a fiber family on
the whole Hermitian parameter space for each prescribed pair of
additive Haar measures.  Thus it applies to every self-dual pair,
not only to one chosen in the proof.  Equivalently one may use the
product construction in AOS10.  It is valid on the independent
locus even at singular parameters.  This verifies the existence
and universal clauses in their original order.  The global measure
compatibility antecedent is retained; as in the original source,
it could also supply these existence clauses directly from its
\texttt{IsOmT} field.
\end{proof}
% END AOS13

% BEGIN AOS14
\paragraph{AOS14. The full original \texttt{HtotInf} for the unchanged stabilizer measure.}
\label{claim:AOS14}
For every original datum under \texttt{HypX},
\[
         \mu_{H,\infty}(H_\infty)=\mu_H(u')(G_\infty),
             \qquad 0<\mu_H(u')(G_\infty)<\infty.
 \tag{AOS.30}
\]
The symbols on the right denote the original ambient real-model
measure and its mass on the whole ambient space, exactly as in
\texttt{HtotInf}; it is supported on its stabilizer.

\begin{proof}
The actual $H_\infty$ is compact by AOS9.  The setting supplies
the Haar-measure property of its unchanged $\mu_{H,\infty}$.
A Haar measure is finite on compact sets and has positive mass
on a nonempty open set, so its total mass is finite and positive.
Now evaluate the second measure tie in (AOS.2) on the whole ambient
group.  Both pushforward maps are measurable, and the preimage of
the whole space is the whole domain.  The resulting equality is
the first equality of (AOS.30).  It transfers both inequalities to
the original $\mu_H(u')$.  This is the full original statement;
the computation does not normalize either stabilizer measure to
probability and does not replace it by the Frobenius comparison
measure used in AOS6--AOS8.
\end{proof}
% END AOS14

% BEGIN AOS15
\paragraph{AOS15. Compact-stabilizer unfolding for every section.}
\label{claim:AOS15}
Let a locally compact second countable group $G$ act on an orbit
$Y=G/H$ with continuous local sections and compact stabilizer $H$.
Let $\mu_G$ be a right invariant left Haar measure, let $\eta$ be
any Haar measure on $H$, and put $\rho=\pi_*\mu_G$.
For every nonnegative measurable $f$ and every section $s$ on the
orbit, even one not assumed measurable,
\[
       \int_G f\,d\mu_G
        =\eta(H)^{-1}
           \int_Y\int_H f(s(y)h)\,d\eta(h)\,d\rho(y).
 \tag{AOS.31}
\]
The same formula holds if $Y$ is regarded as a measurable orbit in
an ambient space and the measures are extended by zero.

\begin{proof}
Compactness and nonzero Haar mass give $0<\eta(H)<\infty$.
Choose a countable cover by continuous section domains and
disjointify it to obtain a Borel section $s_0$.
For this section the inner integral is measurable by Tonelli.
Since $\rho$ is a pushforward, the double integral equals
\[
      \int_G\int_H f(s_0(\pi(g))h)\,d\eta(h)\,d\mu_G(g).
\]
For every $g$, the two representatives $s_0(\pi(g))$ and $g$
differ by right multiplication by an element of $H$.
Left invariance of $\eta$ therefore makes the inner integral
$\int_H f(gh)\,d\eta(h)$.  Interchanging the two nonnegative
integrals and applying right invariance of $\mu_G$ gives
$\eta(H)\int_G f\,d\mu_G$.  Divide by its finite nonzero
factor to obtain (AOS.31), also when the integral of $f$ is infinite.

For any other section $s$, the same left change of variable in $H$
shows pointwise on the orbit that its inner integral equals the
one for $s_0$.  Thus it is the same measurable function there;
no measurable selection property of $s$ is needed.  In an ambient
space the two inner functions agree outside a set of outer measure
zero, so their lower integrals agree as well.  This verifies the
literal all-section form used in the source.
\end{proof}
% END AOS15

% BEGIN AOS16
\paragraph{AOS16. The full original \texttt{WeilInf}.}
\label{claim:AOS16}
Under every original \texttt{HypX}, for all nonnegative measurable
$f$ on the actual $G_\infty$ and all source sections $s$,
\[
 \int_{G_\infty}f(g)\,d\mu_{G,\infty}(g)
   =\mu_{H,\infty}(H_\infty)^{-1}
        \int_{V_\infty}\int_{H_\infty}
                 f(s(y)h)\,d\mu_{H,\infty}(h)
                     \,d(\pi_{Y_0*}\mu_{G,\infty})(y).
 \tag{AOS.32}
\]
This assertion does not require a choice of self-dual measure or
fiber family.

\begin{proof}
The actual $G_\infty$ is a finite product of real unitary groups
through the supplied completed bridge.  It is locally compact,
second countable and unimodular, by the same trace-form argument
used in AOS8.  The stabilizer is the actual $H_\infty$ of AOS9,
with the supplied Haar measure from AOS14.  The orbit map has
continuous local sections: AOS4 gives them at the base tuple,
and translation gives them at every other point.  Its orbit is
the Gram fiber by Witt extension, hence closed because its Gram
matrix is nondegenerate.  The quotient and orbit topologies agree
in these charts.  All hypotheses of AOS15 now hold.  Its formula
is exactly (AOS.32), including the requirement on $s$ only where
an orbit representative exists.  This proves \texttt{WeilInf}
without changing the original $\mu_{H,\infty}$.
\end{proof}
% END AOS16

% BEGIN AOS17
\paragraph{AOS17. The complete original archimedean integration clause.}
\label{claim:AOS17}
For every datum, character and Hermitian witness under exactly the
\texttt{HypX}, \texttt{MeasCompat\,$\psi$}, \texttt{IsStdChar}
and \texttt{GaugeDens} antecedents of AOS11, put
\[
             a_\infty(X,u')
                 =\frac{K_X(u')}{\mu_H(u')(G_\infty)}.
 \tag{AOS.33}
\]
The denominator here is its finite real mass.  This number is
strictly positive, and the full original
$\texttt{AinfClause}(X,\psi,u',h_T,a_\infty)$ holds.  Equivalently,
for every permitted self-dual pair and fiber family, every
nonnegative measurable $f$, and every source section $s$,
\[
       \int f\,d\mu_{G,\infty}
          =a_\infty(X,u')
             \int_{V_\infty}\int_{H_\infty}
                    f(s(y)h)\,d\mu_{H,\infty}(h)\,d\omega_{T_0}(y).
 \tag{AOS.34}
\]

\begin{proof}
AOS13 supplies all existence clauses in the nested \texttt{AllSome}
contract, in its original order.  Fix any witnesses quantified
there.  Substitute (AOS.26) into (AOS.32).  The nonnegative
integral against a scalar multiple of a measure extracts the
scalar $K_X(u')$.  By AOS14 its remaining denominator is exactly
the unchanged finite positive mass $\mu_H(u')(G_\infty)$.
This gives (AOS.34) and the number (AOS.33).  Positivity follows
from positivity of the gauge factors, (AOS.18), $T_w(u')>0$ and
(AOS.30).  For extended nonnegative integrals these operations
remain legitimate because every scalar being divided out is finite
and nonzero.  This proves \texttt{AinfValue} at its full original
scope, including all fiber witnesses and all sections.
\end{proof}
% END AOS17

% BEGIN AOS18
\paragraph{AOS18. The full original \texttt{AprimeInf}: uniqueness of the prescribed scalar.}
\label{claim:AOS18}
For every original datum and every $h_T$ under the antecedents of
AOS11, if $a'>0$ satisfies the original \texttt{AinfClause}, then
\[
       a'=\frac{\operatorname{gaugeU}(X.B.I)
                         \operatorname{volU}(k)^dT_w(u')^{-kd}}
                        {\mu_H(u')(G_\infty)}.
 \tag{AOS.35}
\]
This is the full original \texttt{AprimeInf}, not a statement that
an unspecified positive scalar exists.

\begin{proof}
Choose a compact neighborhood $C$ of the identity in the actual
$G_\infty$.  It has positive finite Haar mass; it is Borel because
the group is Hausdorff.  This construction needs only the supplied
locally compact group and Haar measure.  In particular it gives
the full original \texttt{GinfTest} under \texttt{HypX} alone,
without any character or gauge hypothesis.  Choose one self-dual
pair and one fiber family
whose existence is guaranteed by AOS13, and one Borel section.
Apply the clause for $a'$ and the clause proved in AOS17 to
$f=1_C$.  Both use \emph{the same} witnesses, section and original
stabilizer measure, and hence the same nonnegative inner integral
$I$.  They give
\[
              0<\mu_{G,\infty}(C)=a'I
                      =a_\infty(X,u')I<\infty.
\]
Since $a'$ is a finite positive real number, these equalities
force $0<I<\infty$.  Cancelling $I$ proves
$a'=a_\infty(X,u')$, which is (AOS.35).  The source
\texttt{AprimeInfS} and \texttt{AprimeInf} have this identical
formula and quantifiers, so this proves the literal endpoint.
\end{proof}
% END AOS18

\subsection{Use on the rescaled setting and the remaining consumers}

Apply the chapter to the actual output of GMR19.  Its one standard
character supplies (AOS.5), its preserved measure tie supplies
(AOS.2), and its compatible measure record supplies the original
\texttt{MeasCompat\,$\psi$}.  The original \texttt{HypX} and
\texttt{GaugeDens} antecedents are kept in the statements of the
archimedean consumers; AOS5 supplies the latter for every sign
function.  Every later word is treated using this same
character, bridge and infinity measures.  The auxiliary Frobenius
measures and orthogonal coordinates used to calculate a scalar
have not become fields of a new setting.

The completed endpoint is the exact archimedean integration scalar,
including its all-witness and all-section clauses.  The original
\texttt{RealBracket} still requires its unipotent and local
disintegration factors, and the full mass comparison still requires
the finite-place and arithmetic quotient-volume producers.  In
particular the original \texttt{SiegelWeilUnitaryConvergent}
premise retained by GMR19 is not proved or eliminated from that
upstream original statement by this chapter.  The source handoff
records those next consumers without claiming a whole selected
RH/GRH obligation has been closed.
